% Generated draft: requires mathematical and source review.
\documentclass[11pt]{article}
\usepackage[margin=25mm]{geometry}
\usepackage{fontspec}
\setmainfont{Latin Modern Roman}
\usepackage{amsmath,amssymb,mathtools,unicode-math}
\setmathfont{Latin Modern Math}
\usepackage{enumitem,needspace,xurl,hyperref}
\setlength{\parindent}{0pt}
\setlength{\parskip}{5pt}
\setlength{\emergencystretch}{3em}
\newcommand{\N}{\mathbb{N}}
\newcommand{\Z}{\mathbb{Z}}
\newcommand{\Q}{\mathbb{Q}}
\newcommand{\R}{\mathbb{R}}
\newcommand{\C}{\mathbb{C}}
\newcommand{\F}{\mathbb{F}}
\newcommand{\A}{\mathbb{A}}
\newcommand{\PP}{\mathbb{P}}
\newcommand{\E}{\mathbb{E}}
\newcommand{\eps}{\varepsilon}
\newcommand{\dd}{\,\mathrm{d}}
\begin{document}
\section*{UnsolvedMath ID: 10400100}
\textbf{Source problem code:} catalogue entry
\subsection{Definitions and mathematical statement}
A quandle is a set $X$ with an operation $x\triangleright y$ such that $x\triangleright x=x$, each right translation is bijective, and
\[
(x\triangleright y)\triangleright z=(x\triangleright z)\triangleright(y\triangleright z).
\]
It is connected when the group generated by right translations acts transitively on $X$.

For standard integral quandle homology, start with the free groups $C_n^R(X)=\mathbb Z[X^n]$ and boundary
\[
\partial(x_1,\ldots,x_n)=\sum_{i=2}^n(-1)^i\bigl[(x_1,\ldots,\widehat{x_i},\ldots,x_n)
-(x_1\triangleright x_i,\ldots,x_{i-1}\triangleright x_i,x_{i+1},\ldots,x_n)\bigr].
\]
Quotient by tuples with equal adjacent entries to obtain $C_*^Q(X)$. Then $H_n^Q(X;\mathbb Z)=\ker\partial_n/\operatorname{im}\partial_{n+1}$; with abelian coefficients $A$, cohomology is $H^*(\operatorname{Hom}(C_*^Q(X),A))$.

Compute the degree-two group for every connected quandle and, more generally, develop useful methods to compute quandle homology and cohomology. The flattened notation ``H Q 2'' does not distinguish $H_2^Q(X)$ from $H_Q^2(X)$; the surrounding homology-table reference suggests the former, while the research summary writes the latter. The definitions above distinguish both possibilities without silently identifying them. Integral homology is a usual default, but a cohomology target needs its coefficient convention fixed.

The source does not restrict $X$ to be finite. A general structural formula and an algorithm on finitely encoded quandles are different kinds of methods. Table calculations for examples, Alexander-quandle cohomology results, and Maple programs are supplied partial context, not a computation for every connected quandle.
\subsection{Short English statement}
Compute degree-two quandle (co)homology for every connected quandle and develop broader calculation methods. Resolve the flattened homology/cohomology notation and coefficient convention, and preserve the unrestricted quandle class.
\subsection{Sources}
\begin{itemize}
\item Ohtsuki, \emph{Problems on invariants of knots and 3-manifolds} (2002), Problem 5.4, PDF page 88, \url{https://msp.org/gtm/2002/04/gtm-2002-04-024s.pdf}; accessed 2026-07-29, PDF-font-structure extraction.
\item Source Table 5 and reference [264] for homology examples; [284] for many Alexander-quandle cohomology computations; [185] for Maple programs. The supplied record contains no bibliography entries for these numbered references.
\item Local registry \texttt{data/unsolvedmath/problems.json}, permanent UnsolvedMath ID $10400100$, source code catalogue entry, \emph{Problem 5.4}. Dataset attribution: Ulam AI and the UnsolvedMath Contributors.
\end{itemize}
\end{document}
